This section proves completeness for guarded judgments with arbitrary
qualitative pre- and postconditions.  Let
$P,Q:\Mem\to\Bool$ and $f,g:\Mem\to[0,\infty]$.  For every command $C$ and
post-assertion $(Q\mid g)$, structural induction on $C$ derives the judgment
whose two precomponents are semantically exact: the qualitative component is
$\WP{C}{Q}$ and the quantitative component is $\WP{C}{g}$.  A final use of
\textsc{Consequence} then derives every semantically valid judgment
$\HT{P}{f}{C}{Q}{g}$.

Throughout,
\[
  \bowtie\ \in\ \{\,\leq,=,\geq\,\}
\]
denotes the lower, equality, or upper judgment relation. The induction is
carried out once, for an arbitrary $\bowtie$.

\subsubsection{Exactly what is proved}

For every command $C$, predicates $P,Q$, and expectations $f,g$, we prove
\begin{equation}
  \valid\HT{P}{f}{C}{Q}{g}
  \quad\Longrightarrow\quad
  \vdash\HT{P}{f}{C}{Q}{g}.
  \label{eq:cook-target}
\end{equation}
By the typed validity definition \eqref{eq:cook-validity-typed}, the premise of
\eqref{eq:cook-target} means
\begin{equation}
  \forall m\in\Mem.\quad
  P(m)\Longrightarrow
  \bigl(\WP{C}{Q}(m)\land
        f(m)\bowtie\WP{C}{g}(m)\bigr),
  \label{eq:cook-general-validity}
\end{equation}
or, equivalently,
\begin{equation}
  P\models \WP{C}{Q}\land f\bowtie\WP{C}{g}.
  \label{eq:cook-general-validity-entailment}
\end{equation}
The proof proceeds through the following exact-pre judgment:
\begin{equation}
  \vdash
  \HT{\WP{C}{Q}}{\WP{C}{g}}{C}{Q}{g}
  \label{eq:cook-general-exact-pre-target}
\end{equation}
Here $\WP{C}{Q}$ is the exact qualitative precondition for $Q$, and
$\WP{C}{g}$ is the exact quantitative pre-expectation of $g$.  Once
\eqref{eq:cook-general-exact-pre-target} has been derived, the desired judgment
follows by \textsc{Consequence}:
\begin{equation}
\frac{
  \begin{gathered}
    P\models\WP{C}{Q}\land f\bowtie\WP{C}{g}
    \\[1mm]
    \vdash\HT{\WP{C}{Q}}{\WP{C}{g}}{C}{Q}{g}
    \\[1mm]
    Q\models Q\land g\bowtie g
  \end{gathered}
}{
  \vdash\HT{P}{f}{C}{Q}{g}
}\;\textsc{Consequence}.
\label{eq:cook-general-final-consequence-outline}
\end{equation}
The first premise is \eqref{eq:cook-general-validity-entailment}, the second is
the exact-pre judgment, and the third holds by reflexivity.  It therefore
suffices to prove \eqref{eq:cook-general-exact-pre-target} by structural
induction on $C$.

\subsubsection{The rules and semantic equations used}

The \textsc{Skip}, \textsc{Assign}, \textsc{Sample}, \textsc{Seq}, and
\textsc{If} rules are those in \Cref{fig:hoare:l-u}.  They are parameterised
by $\bowtie$ and therefore cover lower, equality, and upper judgments at once.

\paragraph{The guarded consequence rule.}
Following the actual guarded consequence rule and the notation in the rest of
this note, we use
\begin{equation}
\frac{
  P\models P'\land f\bowtie f'
  \qquad
  \vdash\HT{P'}{f'}{C}{Q'}{g'}
  \qquad
  Q'\models Q\land g'\bowtie g
}{
  \vdash\HT{P}{f}{C}{Q}{g}
}\;\textsc{Consequence}.
\label{eq:cook-consequence}
\end{equation}
For $\bowtie={\leq}$ this is the lower rule, for $\bowtie={=}$ the equality
rule, and for $\bowtie={\geq}$ the upper rule.  The first comparison is
required on the new pre-guard $P$, and the last comparison is required on the
source post-guard $Q'$.  Every use of \textsc{Consequence} below is displayed
in this three-premise form with all entries already instantiated.

\paragraph{Semantic equations.}
The expectation semantics gives
\begin{align}
  \WP{\SKIP}{g} &= g,
  \label{eq:cook-wp-skip}\\
  \WP{\vx <- \expr}{g} &= g[\vx\mapped\expr],
  \label{eq:cook-wp-assign}\\
  \WP{\vx <* \sexpr}{g}
    &= \E{\sexpr}{\lambda\val.\,g[\vx\mapped\val]},
  \label{eq:cook-wp-sample}\\
  \WP{C_1;C_2}{g}
    &= \WP{C_1}{\WP{C_2}{g}},
  \label{eq:cook-wp-seq}\\
  \WP{\ITE{B}{C_1}{C_2}}{g}
    &= [B]\cdot\WP{C_1}{g}
       +[\neg B]\cdot\WP{C_2}{g}.
  \label{eq:cook-wp-if}
\end{align}
These equations follow from the denotational clauses above rather than being
additional proof rules: the first four follow from the return and bind laws,
and the conditional equation follows by splitting on $B(m)$.

For $L=\WHILE B\DO D$, define
\begin{equation}
  \Phi_g(X)\defsym[B]\cdot\WP{D}{X}+[\neg B]\cdot g.
  \label{eq:cook-loop-functional}
\end{equation}
The loop semantics used in this note is
\begin{equation}
  \WP{L}{g}=\operatorname{lfp}(\Phi_g)
  =\sup_{n\geq0}\Phi_g^n(0).
  \label{eq:cook-loop-lfp}
\end{equation}

We assume that the exact preconditions, pre-expectations, and loop residuals
used below are expressible, and that the background assertion theory proves
the required semantic equalities and inequalities.  The loop case additionally
uses monotonicity, Scott continuity, and additivity of $\mathsf{wp}$.  Its exact
pre-expectation is assumed finite-valued so that residual subtraction never
forms $\infty-\infty$.

\subsubsection{The exact-pre lemma}

\begin{lemma}[Exact-pre judgments]
\label{lem:exact-pre-derivable}
For every command $C$, every postcondition $Q$, every post-expectation $g$ in
the expressive, finitary fragment, and each
$\bowtie\in\{\leq,=,\geq\}$,
\begin{equation}
  \vdash\HT{\WP{C}{Q}}{\WP{C}{g}}{C}{Q}{g}.
  \label{eq:cook-exact-bowtie}
\end{equation}
\end{lemma}

\begin{proof}
The proof is by structural induction on $C$, with $Q$, $g$, and $\bowtie$
universally quantified.

\paragraph{Case: skip.}

The typed weakest-precondition equations are
\[
  \WP{\SKIP}{Q}=Q,
  \qquad
  \WP{\SKIP}{g}=g.
\]
Hence \textsc{Skip} gives
\[
\frac{}{\vdash\HT{Q}{g}{\SKIP}{Q}{g}}\;\textsc{Skip},
\]
which is \eqref{eq:cook-exact-bowtie} after rewriting both precomponents.

\paragraph{Case: assignment.}

The typed assignment equations are
\[
  \WP{\vx <- \expr}{Q}=Q[\vx\mapped\expr],
  \qquad
  \WP{\vx <- \expr}{g}=g[\vx\mapped\expr].
\]
The required instance of \textsc{Assign} is
\[
\frac{}{
  \vdash\HT{Q[\vx\mapped\expr]}{g[\vx\mapped\expr]}
    {\vx <- \expr}{Q}{g}
}\;\textsc{Assign}.
\]
Rewriting its two precomponents proves \eqref{eq:cook-exact-bowtie}.

\paragraph{Case: sampling.}

The typed sampling equations are
\[
  \WP{\vx <* \sexpr}{Q}
  =\forall\val\in\supp(\sexpr).\ Q[\vx\mapped\val],
  \qquad
  \WP{\vx <* \sexpr}{g}
  =\E{\sexpr}{\lambda\val.\,g[\vx\mapped\val]}.
\]
Thus \textsc{Sample} gives
\[
\frac{}{
  \vdash\HT{\forall\val\in\supp(\sexpr).\ Q[\vx\mapped\val]}
    {\E{\sexpr}{\lambda\val.\,g[\vx\mapped\val]}}
    {\vx <* \sexpr}{Q}{g}
}\;\textsc{Sample}.
\]
Rewriting its two precomponents proves \eqref{eq:cook-exact-bowtie}.

\paragraph{Case: sequential composition.}

Define
\[
  R\defsym\WP{C_2}{Q},
  \qquad
  h\defsym\WP{C_2}{g}.
\]
The two induction hypotheses are
\[
  \vdash\HT{R}{h}{C_2}{Q}{g},
  \qquad
  \vdash\HT{\WP{C_1}{R}}{\WP{C_1}{h}}{C_1}{R}{h}.
\]
They are exactly the two premises of \textsc{Seq}:
\begin{equation}
\frac{
  \vdash\HT{\WP{C_1}{R}}{\WP{C_1}{h}}{C_1}{R}{h}
  \qquad
  \vdash\HT{R}{h}{C_2}{Q}{g}
}{
  \vdash\HT{\WP{C_1}{R}}{\WP{C_1}{h}}{C_1;C_2}{Q}{g}
}\;\textsc{Seq}.
\label{eq:cook-seq-derivation}
\end{equation}
The composition equations
\[
  \WP{C_1;C_2}{Q}=\WP{C_1}{R},
  \qquad
  \WP{C_1;C_2}{g}=\WP{C_1}{h}
\]
rewrite the conclusion to \eqref{eq:cook-exact-bowtie}.

\paragraph{Case: $C$ is a conditional on $B$.}

The structurally smaller commands are $C_1$ and $C_2$.  Fix $g$ and define
\begin{equation}
  H\defsym[B]\cdot\WP{C_1}{g}
    +[\neg B]\cdot\WP{C_2}{g}.
  \label{eq:cook-if-H}
\end{equation}
By \eqref{eq:cook-wp-if},
\begin{equation}
  H=\WP{\ITE{B}{C_1}{C_2}}{g}.
  \label{eq:cook-if-H-wp}
\end{equation}
The indicator semantics gives
\begin{align}
  B&\models H=\WP{C_1}{g},
  \label{eq:cook-if-true}\\
  \neg B&\models H=\WP{C_2}{g}.
  \label{eq:cook-if-false}
\end{align}
These are semantic facts, not induction hypotheses or syntactic judgments.

For the arbitrary relation $\bowtie$, the two IHs are
\begin{align}
  &\vdash\HT{\top}{\WP{C_1}{g}}{C_1}{\top}{g},
  \label{eq:cook-if-IH-one}\\
  &\vdash\HT{\top}{\WP{C_2}{g}}{C_2}{\top}{g}.
  \label{eq:cook-if-IH-two}
\end{align}

\paragraph{The true branch by \textsc{Consequence}.}
Substitute
\[
  P=B,\quad f=H,\quad P'=\top,\quad f'=\WP{C_1}{g},
  \quad Q'=Q=\top,\quad g'=g
\]
in \eqref{eq:cook-consequence}.  The complete rule instance is
\begin{equation}
\frac{
  \begin{array}{c}
    B\models\top\land H\bowtie\WP{C_1}{g}
    \qquad
    \vdash\HT{\top}{\WP{C_1}{g}}{C_1}{\top}{g}
    \\[1mm]
    \top\models\top\land g\bowtie g
  \end{array}
}{
  \vdash\HT{B}{H}{C_1}{\top}{g}
}\;\textsc{Consequence}.
\label{eq:cook-if-conseq-true}
\end{equation}
The middle premise is IH~\eqref{eq:cook-if-IH-one}.  The first premise holds
because $B\models\top$ and \eqref{eq:cook-if-true} is an equality, hence it
implies each of $H\leq\WP{C_1}{g}$, $H=\WP{C_1}{g}$, and
$H\geq\WP{C_1}{g}$.  The last premise holds by reflexivity.  Thus the same
displayed derivation is valid for all three choices of $\bowtie$.

\paragraph{The false branch by \textsc{Consequence}.}
Now substitute
\[
  P=\neg B,\quad f=H,\quad P'=\top,
  \quad f'=\WP{C_2}{g},
  \quad Q'=Q=\top,\quad g'=g.
\]
The complete rule instance is
\begin{equation}
\frac{
  \begin{array}{c}
    \neg B\models\top\land H\bowtie\WP{C_2}{g}
    \qquad
    \vdash\HT{\top}{\WP{C_2}{g}}{C_2}{\top}{g}
    \\[1mm]
    \top\models\top\land g\bowtie g
  \end{array}
}{
  \vdash\HT{\neg B}{H}{C_2}{\top}{g}
}\;\textsc{Consequence}.
\label{eq:cook-if-conseq-false}
\end{equation}
The middle premise is IH~\eqref{eq:cook-if-IH-two}; the first follows from
the semantic equality \eqref{eq:cook-if-false}; and the last is reflexive.

\paragraph{Application of \textsc{If}.}
The conclusions of \eqref{eq:cook-if-conseq-true} and
\eqref{eq:cook-if-conseq-false} fill the two branch premises of
\textsc{If}:
\begin{equation}
\frac{
  \vdash\HT{B}{H}{C_1}{\top}{g}
  \qquad
  \vdash\HT{\neg B}{H}{C_2}{\top}{g}
}{
  \vdash\HT{\top}{H}{\ITE{B}{C_1}{C_2}}{\top}{g}
}\;\textsc{If}.
\label{eq:cook-if-derivation}
\end{equation}
The branch guards are $B\land\top=B$ and
$\neg B\land\top=\neg B$.  Finally,
\eqref{eq:cook-if-H-wp} rewrites $H$ to the exact pre-expectation.  This
proves the conditional case simultaneously for
$\bowtie\in\{\leq,=,\geq\}$.

\paragraph{Case: $C$ is a loop $L$ with guard $B$ and body $D$.}

\noindent\todo{TODO (equality case for while): Add an
equality-introduction rule combining the exact lower and upper loop judgments,
or provide a dedicated equality while rule. Until then, the derivations below
establish the two inequalities but do not yet construct the syntactic $[=]$
judgment.}\par

The only structurally smaller command is the body $D$.  Fix $g$ and define
\[
  W\defsym\WP{L}{g}.
\]
This is the semantic loop invariant: it is the exact expected value of $g$
after the whole loop.

\paragraph{Semantic fixed-point facts.}
Let
\[
  K_0\defsym0,
  \qquad
  K_{n+1}\defsym\Phi_g(K_n).
\]
By \eqref{eq:cook-loop-lfp} and Scott continuity,
\begin{equation}
  K_n\uparrow W,
  \qquad
  W=\Phi_g(W).
  \label{eq:cook-loop-kleene}
\end{equation}
Unfolding $\Phi_g$ and evaluating indicators gives the semantic equalities
\begin{align}
  B&\models W=\WP{D}{W},
  \label{eq:cook-loop-on-B}\\
  \neg B&\models W=g,
  \label{eq:cook-loop-on-not-B}\\
  B&\models K_{n+1}=\WP{D}{K_n}\qquad(n\geq0).
  \label{eq:cook-approximant-on-B}
\end{align}
None of these equalities is itself a syntactic Hoare derivation.

\paragraph{The body judgment, once for all three relations.}
The IH for $D$, instantiated at post-expectation $W$ and arbitrary relation
$\bowtie$, is
\begin{equation}
  \vdash\HT{\top}{\WP{D}{W}}{D}{\top}{W}.
  \label{eq:cook-loop-body-IH}
\end{equation}
Substitute
\[
  P=B,\quad f=W,\quad P'=\top,\quad f'=\WP{D}{W},
  \quad Q'=Q=\top,\quad g'=W
\]
in \textsc{Consequence}.  This gives the full rule instance
\begin{equation}
\frac{
  \begin{array}{c}
    B\models\top\land W\bowtie\WP{D}{W}
    \qquad
    \vdash\HT{\top}{\WP{D}{W}}{D}{\top}{W}
    \\[1mm]
    \top\models\top\land W\bowtie W
  \end{array}
}{
  \vdash\HT{B}{W}{D}{\top}{W}
}\;\textsc{Consequence}.
\label{eq:cook-loop-body-bowtie}
\end{equation}
The first premise follows from equality \eqref{eq:cook-loop-on-B}, the middle
premise is IH~\eqref{eq:cook-loop-body-IH}, and the last premise is reflexive.
The lower instance of \eqref{eq:cook-loop-body-bowtie} will be Premise~1 of
\textsc{While-lb}; its upper instance will be the body premise of
\textsc{While-ub}.

\paragraph{Residual witnesses for the lower loop rule.}
Because $K_n\uparrow W$, we have $0\leq K_n\leq W$.  Finiteness of $W$
allows us to define
\begin{equation}
  r_n\defsym W-K_n.
  \label{eq:cook-residual-definition}
\end{equation}
Each $r_n$ is a non-negative, finite-valued expectation.

\paragraph{Lower Premise 1.}
Instantiate \textsc{While-lb} with qualitative loop invariant $P:=\top$ and
quantitative loop invariant $W$.  Its first premise is
\[
  \vdash\HT[lower]{B}{W}{D}{\top}{W},
\]
which is exactly the $\bowtie={\leq}$ instance of
\eqref{eq:cook-loop-body-bowtie}. 

\paragraph{Lower Premise 2.}
The second premise is the semantic side condition
\[
  B\models r_0\geq W.
\]
Since $K_0=0$, definition \eqref{eq:cook-residual-definition} gives
$r_0=W-K_0=W$.  Hence $B\models r_0=W\geq W$.

\paragraph{Semantic residual equation for Lower Premise 3.}
For every $n$,
\begin{equation}
  B\models r_{n+1}=\WP{D}{r_n}.
  \label{eq:cook-residual-step}
\end{equation}
Indeed, $W=K_n+r_n$.  Additivity of $\mathsf{wp}$ gives
\[
  \WP{D}{W}=\WP{D}{K_n}+\WP{D}{r_n}.
\]
On $B$, equations \eqref{eq:cook-loop-on-B} and
\eqref{eq:cook-approximant-on-B} turn this into
\[
  W=K_{n+1}+\WP{D}{r_n}.
\]
Definition \eqref{eq:cook-residual-definition} also gives
$W=K_{n+1}+r_{n+1}$.  All terms are finite-valued, so cancellation yields
\eqref{eq:cook-residual-step}.  The equality appears because the chosen
$r_n=W-K_n$ is the \emph{exact} residual, not merely an arbitrary upper bound.
Premise~3 of \textsc{While-lb} only needs
$B\models r_{n+1}\geq\WP{D}{r_n}$, so
\eqref{eq:cook-residual-step} is a strictly stronger semantic fact.  It is not
yet a syntactic judgment; the required upper judgment is derived next.

\paragraph{Lower Premise 3.}
Fix an arbitrary $n$.  The upper IH for $D$, instantiated at post-expectation
$r_n$, is
\begin{equation}
  \vdash\HT[upper]{\top}{\WP{D}{r_n}}{D}{\top}{r_n}.
  \label{eq:cook-loop-residual-upper-IH}
\end{equation}
Substituting
\[
  P=B,\quad f=r_{n+1},\quad P'=\top,
  \quad f'=\WP{D}{r_n},
  \quad Q'=Q=\top,\quad g'=r_n,
  \quad\bowtie={\geq}
\]
in \textsc{Consequence} gives
\begin{equation}
\frac{
  \begin{array}{c}
    B\models\top\land r_{n+1}\geq\WP{D}{r_n}
    \qquad
    \vdash\HT[upper]{\top}{\WP{D}{r_n}}{D}{\top}{r_n}
    \\[1mm]
    \top\models\top\land r_n\geq r_n
  \end{array}
}{
  \vdash\HT[upper]{B}{r_{n+1}}{D}{\top}{r_n}
}\;\textsc{Consequence}.
\label{eq:cook-loop-residual-conseq}
\end{equation}
The first premise follows from \eqref{eq:cook-residual-step}, the middle one
is IH~\eqref{eq:cook-loop-residual-upper-IH}, and the last is reflexive.
This use of \textsc{Consequence} is essential: the IH has pre-guard $\top$ and
pre-expectation $\WP{D}{r_n}$, whereas the premise required by
\textsc{While-lb} has pre-guard $B$ and pre-expectation $r_{n+1}$.  The first
side condition performs exactly these two adaptations under the assumption
$B$.
Because $n$ was arbitrary, this derives
\[
  \forall n\geq0.\quad
  \vdash\HT[upper]{B}{r_{n+1}}{D}{\top}{r_n}.
\]

\paragraph{Lower Premise 4.}
The fourth premise is
\[
  B\models\lim_{n\to\infty}r_n=0.
\]
By \eqref{eq:cook-loop-kleene}, $K_n\uparrow W$.  Since $W$ is finite-valued,
pointwise
\[
  \lim_{n\to\infty}r_n
  =\lim_{n\to\infty}(W-K_n)
  =W-\lim_{n\to\infty}K_n
  =W-W=0.
\]
The equality holds on all states and thus on $B$.

\paragraph{Application of \textsc{While-lb}.}
The four established premises fill the lower loop rule as follows:
\begin{equation}
\frac{
  \begin{array}{c}
    \vdash\HT[lower]{B}{W}{D}{\top}{W}
    \qquad
    B\models r_0\geq W
    \\[1mm]
    \forall n.\ \vdash\HT[upper]{B}{r_{n+1}}{D}{\top}{r_n}
    \qquad
    B\models\lim_{n\to\infty}r_n=0
  \end{array}
}{
  \vdash\HT[lower]{\top}{W}{L}{\neg B}{W}
}\;\textsc{While-lb}.
\label{eq:cook-loop-lower-canonical}
\end{equation}

\paragraph{Application of \textsc{While-ub}.}
The $\bowtie={\geq}$ instance of \eqref{eq:cook-loop-body-bowtie} is
\[
  \vdash\HT[upper]{B}{W}{D}{\top}{W}.
\]
It is exactly the body premise obtained by taking $P=\top$ and $v=W$ in the
canonical upper loop rule:
\begin{equation}
\frac{
  \vdash\HT[upper]{B}{W}{D}{\top}{W}
}{
  \vdash\HT[upper]{\top}{W}{L}{\neg B}{W}
}\;\textsc{While-ub}.
\label{eq:cook-loop-upper-canonical}
\end{equation}

\paragraph{The common final consequence step.}
Equations \eqref{eq:cook-loop-lower-canonical} and
\eqref{eq:cook-loop-upper-canonical} establish, for the corresponding choice
of $\bowtie$,
\[
  \vdash\HT{\top}{W}{L}{\neg B}{W}.
\]
Both are converted to the desired post-assertion by the same fully
instantiated consequence rule:
\begin{equation}
\frac{
  \top\models\top\land W\bowtie W
  \qquad
  \vdash\HT{\top}{W}{L}{\neg B}{W}
  \qquad
  \neg B\models\top\land W\bowtie g
}{
  \vdash\HT{\top}{W}{L}{\top}{g}
}\;\textsc{Consequence}.
\label{eq:cook-loop-final-conseq}
\end{equation}
The first premise is reflexive.  The middle premise comes from
\textsc{While-lb} when $\bowtie={\leq}$ and from \textsc{While-ub} when
$\bowtie={\geq}$.  The last premise follows from
$\neg B\models W=g$ in \eqref{eq:cook-loop-on-not-B}; equality implies both
directions of $\bowtie$.  Since $W=\WP{L}{g}$ by definition, the conclusion
of \eqref{eq:cook-loop-final-conseq} is the exact-pre judgment for the loop.
This closes the lower and upper branches of the structural induction; the
equality branch is the explicit TODO recorded at the start of this case.
\end{proof}

\subsubsection{From the exact pre-expectation to every valid lower bound}

\begin{theorem}[Completeness for lower bounds]
\label{thm:cook-relative-completeness}
For every command $C$ and expectations $f,g$ in the expressive, finitary
fragment,
\[
  \valid\HT[lower]{\top}{f}{C}{\top}{g}
  \quad\Longrightarrow\quad
  \vdash\HT[lower]{\top}{f}{C}{\top}{g}.
\]
\end{theorem}

\begin{proof}
Assume
\[
  \valid\HT[lower]{\top}{f}{C}{\top}{g}.
\]
By semantic validity, the qualitative obligation
$\supp(C_m)\subseteq\top$ is automatic, while the quantitative obligation is
\[
  \forall m.\quad
  f(m)\leq\E{C,m}{[\top]\cdot g}.
\]
Because $[\top]=1$, definition \eqref{eq:cook-wp-definition} reduces it to
\begin{equation}
  f\leq\WP{C}{g}.
  \label{eq:cook-valid-bound}
\end{equation}
The $\bowtie={\leq}$ instance of the exact-pre lemma gives
\begin{equation}
  \vdash\HT[lower]{\top}{\WP{C}{g}}{C}{\top}{g}.
  \label{eq:cook-final-exact}
\end{equation}
Now substitute
\[
  P=P'=Q'=Q=\top,
  \quad f'=\WP{C}{g},
  \quad g'=g,
  \quad\bowtie={\leq}
\]
in \textsc{Consequence}.  The complete final rule instance is
\begin{equation}
\frac{
  \top\models\top\land f\leq\WP{C}{g}
  \qquad
  \vdash\HT[lower]{\top}{\WP{C}{g}}{C}{\top}{g}
  \qquad
  \top\models\top\land g\leq g
}{
  \vdash\HT[lower]{\top}{f}{C}{\top}{g}
}\;\textsc{Consequence}.
\label{eq:cook-final-consequence}
\end{equation}
The first premise is \eqref{eq:cook-valid-bound}, the middle premise is
\eqref{eq:cook-final-exact}, and the last premise is reflexive.  The
conclusion is exactly the desired derivability statement.
\end{proof}